\section{Scope and Open Problems}

A theory is honest in proportion to how clearly it marks its own edges, so this section says what the paper covers, what it reaches toward without quite grasping, and what it leaves alone, and then names the genuine open problems plainly. None of these are hidden, and the list moves, but the ones that remain are all at the level of physics.

The scope falls into three honest bands. Some results are established, derived by exact computation or measured on the substrate with a control, and these are the solid rows of the scoreboard, the spin and the fermion, the Dirac cone and the $g$-factor, the gauge structure and the unification ratios, the holography, the selection of the substrate, the self that binds itself. Some are extrapolated, grounded in the structure but not yet shown end to end, and these include the gravitational and cosmological results at the effective level and the whole inner reading, which rests throughout on the one premise. And some are deliberately out of scope, the absolute low-energy spectrum, a fully dynamical derivation of the inner architecture, anything that would require the exact collision run at scales beyond reach. The paper claims the first band, reaches for the second with the seam marked, and does not pretend to the third.

\subsection{The open frontiers}

The central open problem, behind several others, is to feed the actual four-dimensional collision, not a reduced toy, through the renormalization blocking that is already a checked method, and read off the absolute Standard Model coefficients. The mechanism of the running is in hand, the make-or-break simulation is not, and behind it sit all the absolute numbers, the mass spectrum, the gauge-coupling values, the mixing angles, none of which the model yet computes. Closely related is the value of the couplings and masses, where the bare rule fixes every form but no magnitude, and deriving a number like the fine-structure constant needs a further principle the rule does not contain. The generation count is solid but its splitting into three distinct masses is the open conjecture noted earlier. Curved-bulk gravity on $\{5,3,4\}$ did not yield a clean discriminator at reachable sizes and is reported as open rather than forced.

Two further frontiers concern how matter and selves assemble. Matter does not yet condense onto the cusp from a generic start, the bare knit from equilibrium churning rather than concentrating, exactly as a single reversible conserving rule must from an equilibrium start, so clean settling needs either special low-entropy initial conditions, which the growing edge supplies, or a deeper demonstration. And while a single self binds itself, clean binding between separated selves, the cross-gap attraction that would build the higher rungs of the tower as a dynamical force rather than a coarse-grained pattern, is not yet measured. The full continuum limit, the restoration of the continuous symmetries from the discrete base across the whole tower at once, is established in pieces and open as one chain.

One question is open on purpose and left to the reader. Whether the universe has a true first beat or has been growing from beyond any beginning is not settled by anything in the law, and both readings are consistent with it. We let it stand, and meet it once more at the close. The frontiers above are the dark parts of the fire, the next places to bring light, and naming them is the price of being trusted with the parts that already burn.
